% Options for packages loaded elsewhere
\PassOptionsToPackage{unicode}{hyperref}
\PassOptionsToPackage{hyphens}{url}
\PassOptionsToPackage{dvipsnames,svgnames,x11names}{xcolor}
\documentclass[
  spanish,
  12pt,
  a4paper,
]{article}
\usepackage{xcolor}
\usepackage[margin=2.54cm]{geometry}
\usepackage{amsmath,amssymb}
\setcounter{secnumdepth}{-\maxdimen} % remove section numbering
\usepackage{iftex}
\ifPDFTeX
  \usepackage[T1]{fontenc}
  \usepackage[utf8]{inputenc}
  \usepackage{textcomp} % provide euro and other symbols
\else % if luatex or xetex
  \usepackage{unicode-math} % this also loads fontspec
  \defaultfontfeatures{Scale=MatchLowercase}
  \defaultfontfeatures[\rmfamily]{Ligatures=TeX,Scale=1}
\fi
\usepackage{lmodern}
\ifPDFTeX\else
  % xetex/luatex font selection
  \setmainfont[]{Times New Roman}
\fi
% Use upquote if available, for straight quotes in verbatim environments
\IfFileExists{upquote.sty}{\usepackage{upquote}}{}
\IfFileExists{microtype.sty}{% use microtype if available
  \usepackage[]{microtype}
  \UseMicrotypeSet[protrusion]{basicmath} % disable protrusion for tt fonts
}{}
\usepackage{setspace}
\makeatletter
\@ifundefined{KOMAClassName}{% if non-KOMA class
  \IfFileExists{parskip.sty}{%
    \usepackage{parskip}
  }{% else
    \setlength{\parindent}{0pt}
    \setlength{\parskip}{6pt plus 2pt minus 1pt}}
}{% if KOMA class
  \KOMAoptions{parskip=half}}
\makeatother
\usepackage{color}
\usepackage{fancyvrb}
\newcommand{\VerbBar}{|}
\newcommand{\VERB}{\Verb[commandchars=\\\{\}]}
\DefineVerbatimEnvironment{Highlighting}{Verbatim}{commandchars=\\\{\}}
% Add ',fontsize=\small' for more characters per line
\newenvironment{Shaded}{}{}
\newcommand{\AlertTok}[1]{\textcolor[rgb]{1.00,0.00,0.00}{\textbf{#1}}}
\newcommand{\AnnotationTok}[1]{\textcolor[rgb]{0.38,0.63,0.69}{\textbf{\textit{#1}}}}
\newcommand{\AttributeTok}[1]{\textcolor[rgb]{0.49,0.56,0.16}{#1}}
\newcommand{\BaseNTok}[1]{\textcolor[rgb]{0.25,0.63,0.44}{#1}}
\newcommand{\BuiltInTok}[1]{\textcolor[rgb]{0.00,0.50,0.00}{#1}}
\newcommand{\CharTok}[1]{\textcolor[rgb]{0.25,0.44,0.63}{#1}}
\newcommand{\CommentTok}[1]{\textcolor[rgb]{0.38,0.63,0.69}{\textit{#1}}}
\newcommand{\CommentVarTok}[1]{\textcolor[rgb]{0.38,0.63,0.69}{\textbf{\textit{#1}}}}
\newcommand{\ConstantTok}[1]{\textcolor[rgb]{0.53,0.00,0.00}{#1}}
\newcommand{\ControlFlowTok}[1]{\textcolor[rgb]{0.00,0.44,0.13}{\textbf{#1}}}
\newcommand{\DataTypeTok}[1]{\textcolor[rgb]{0.56,0.13,0.00}{#1}}
\newcommand{\DecValTok}[1]{\textcolor[rgb]{0.25,0.63,0.44}{#1}}
\newcommand{\DocumentationTok}[1]{\textcolor[rgb]{0.73,0.13,0.13}{\textit{#1}}}
\newcommand{\ErrorTok}[1]{\textcolor[rgb]{1.00,0.00,0.00}{\textbf{#1}}}
\newcommand{\ExtensionTok}[1]{#1}
\newcommand{\FloatTok}[1]{\textcolor[rgb]{0.25,0.63,0.44}{#1}}
\newcommand{\FunctionTok}[1]{\textcolor[rgb]{0.02,0.16,0.49}{#1}}
\newcommand{\ImportTok}[1]{\textcolor[rgb]{0.00,0.50,0.00}{\textbf{#1}}}
\newcommand{\InformationTok}[1]{\textcolor[rgb]{0.38,0.63,0.69}{\textbf{\textit{#1}}}}
\newcommand{\KeywordTok}[1]{\textcolor[rgb]{0.00,0.44,0.13}{\textbf{#1}}}
\newcommand{\NormalTok}[1]{#1}
\newcommand{\OperatorTok}[1]{\textcolor[rgb]{0.40,0.40,0.40}{#1}}
\newcommand{\OtherTok}[1]{\textcolor[rgb]{0.00,0.44,0.13}{#1}}
\newcommand{\PreprocessorTok}[1]{\textcolor[rgb]{0.74,0.48,0.00}{#1}}
\newcommand{\RegionMarkerTok}[1]{#1}
\newcommand{\SpecialCharTok}[1]{\textcolor[rgb]{0.25,0.44,0.63}{#1}}
\newcommand{\SpecialStringTok}[1]{\textcolor[rgb]{0.73,0.40,0.53}{#1}}
\newcommand{\StringTok}[1]{\textcolor[rgb]{0.25,0.44,0.63}{#1}}
\newcommand{\VariableTok}[1]{\textcolor[rgb]{0.10,0.09,0.49}{#1}}
\newcommand{\VerbatimStringTok}[1]{\textcolor[rgb]{0.25,0.44,0.63}{#1}}
\newcommand{\WarningTok}[1]{\textcolor[rgb]{0.38,0.63,0.69}{\textbf{\textit{#1}}}}
\usepackage{longtable,booktabs,array}
\usepackage{caption}
\captionsetup[table]{skip=6pt}
\newcounter{none} % for unnumbered tables
\usepackage{calc} % for calculating minipage widths
% Correct order of tables after \paragraph or \subparagraph
\usepackage{etoolbox}
\makeatletter
\patchcmd\longtable{\par}{\if@noskipsec\mbox{}\fi\par}{}{}
\makeatother
\ifLuaTeX
\usepackage[bidi=basic,shorthands=off]{babel}
\else
\usepackage[bidi=default,shorthands=off]{babel}
\fi
\ifPDFTeX
\else
\babelfont{rm}[]{Times New Roman}
\fi
\ifLuaTeX
  \usepackage{selnolig} % disable illegal ligatures
\fi
\setlength{\emergencystretch}{3em} % prevent overfull lines
\providecommand{\tightlist}{%
  \setlength{\itemsep}{0pt}\setlength{\parskip}{0pt}}
\usepackage{hanging}
\usepackage{bookmark}
\IfFileExists{xurl.sty}{\usepackage{xurl}}{} % add URL line breaks if available
\urlstyle{same}
% fallback for those not using the hyperref driver hyperxmp:
\makeatletter
\@ifundefined{xmpquote}{\newcommand{\xmpquote}[1]{#1}}{}
\makeatother
\hypersetup{
  pdftitle={Sistema de gestión de residuos y reciclaje incentivado para comercios y centros urbanos},
  pdfauthor={Equipo del proyecto},
  pdflang={es-SV},
  colorlinks=true,
  linkcolor={blue},
  filecolor={Maroon},
  citecolor={Blue},
  urlcolor={Blue},
  pdfcreator={LaTeX via pandoc}}

\title{Sistema de gestión de residuos y reciclaje incentivado para
comercios y centros urbanos}
\usepackage{etoolbox}
\makeatletter
\providecommand{\subtitle}[1]{% add subtitle to \maketitle
  \apptocmd{\@title}{\par {\large #1 \par}}{}{}
}
\makeatother
\subtitle{Memoria técnica del Proyecto 15}
\author{Equipo del proyecto}
\date{23 de noviembre de 2026}

\begin{document}
\maketitle

{
\setcounter{tocdepth}{2}
\tableofcontents
}
\setstretch{2}
\newpage

\section{1. Portada}\label{portada}

\textbf{Sistema de gestión de residuos y reciclaje incentivado para
comercios y centros urbanos}\\
Memoria técnica del Proyecto 15\\
Fecha de corte documental: 23 de noviembre de 2026\\
Datos institucionales y de autoría individual: no constan en el material
del repositorio.

\textbf{Nota de alcance.} Esta memoria consolida la evidencia existente
en el repositorio. No se localizó una memoria principal anterior en
Markdown ni un PDF final para actualizar; por eso este archivo establece
la versión consolidada actual. Los datos estadísticos se transcriben de
\texttt{results/analysis.json} y las observaciones de
\texttt{results/observations.csv}. Las limitaciones de evidencia se
declaran en cada fase.

\section{2. Resumen}\label{resumen}

Se implementó un sistema web para gestionar solicitudes de recolección
de residuos de comercios, asignación de recolectores, registro y
verificación de pesajes, acreditación de puntos y canje de recompensas.
El backend está desarrollado en Rust y persiste documentos en Apache
CouchDB. El cliente TypeScript usa PouchDB para conservar datos y
cambios locales en el navegador y sincronizarlos con la API.

La evaluación de carga compara condiciones de red del mismo API
Rust+CouchDB, no motores de base de datos ni APIs completas
equivalentes. \texttt{analysis.json} incluye tres corridas válidas por
condición: ideal (media 38.9878 solicitudes/s; DE 4.1772) e intermitente
(media 9.2737 solicitudes/s; DE 0.5474). La prueba global fue ANOVA,
\emph{F} = 149.2368, \emph{p} = 0.0002578. La comparación post hoc Welch
tuvo \emph{p} ajustado de Holm = 0.0058430 y delta de Cliff = 1.0. La
muestra inferencial contiene únicamente corridas con flujo de negocio
completamente exitoso y cero fallos HTTP. Los resultados representan el
entorno de laboratorio y no demuestran por sí solos una ventaja de
producto ni permiten atribuir causalidad a un componente aislado.

El baseline PostgreSQL/\texttt{pgbench} mide upsert y lectura indexada
sobre una tabla, no un servidor API alternativo. Las series de CPU y RAM
atribuibles al contenedor API están vacías en las corridas archivadas y
los hilos no fueron medidos. La revisión bibliográfica fue exploratoria:
la matriz registra antecedentes preseleccionados, pero no un cribado
sistemático terminado. Los datos locales web se almacenan en
IndexedDB/PouchDB y el JWT en \texttt{localStorage}; la aplicación no
cifra estos datos en reposo.

\textbf{Palabras clave:} reciclaje, comercio, incentivos, Rust, CouchDB,
PouchDB, sincronización offline, red degradada.

\section{3. Introducción}\label{introducciuxf3n}

La gestión de materiales reciclables requiere coordinar generadores,
recolectores y centros de acopio, conservar evidencia de las cantidades
recibidas y vincular esa actividad con incentivos verificables. El
Proyecto 15 aborda esa coordinación por medio de una aplicación web con
flujo de trabajo por roles y operación tolerante a interrupciones de
conectividad.

La memoria presenta el problema, el respaldo bibliográfico disponible,
el diseño y la implementación, y los resultados experimentales
efectivamente registrados. Distingue entre funcionalidades
implementadas, herramientas listas para ejecutarse y resultados ya
medidos. Esta separación evita convertir una configuración experimental
en un resultado o una búsqueda exploratoria en una revisión sistemática.

\section{4. Planteamiento del problema y
contexto}\label{planteamiento-del-problema-y-contexto}

El problema de diseño es cómo registrar y coordinar solicitudes de
recolección, pesajes y recompensas cuando los participantes cumplen
funciones distintas y la conexión puede ser intermitente. El sistema
debe preservar trazabilidad de las transiciones y evitar duplicar
pesajes o puntos al reintentar una sincronización.

El repositorio incorpora como fuente de contexto local un diagnóstico
institucional del mercado de reciclaje de El Salvador (Ministerio de
Medio Ambiente y Recursos Naturales {[}MARN{]}, 2023). La matriz
bibliográfica informa que se revisó la ficha y no el informe completo;
por ello esta memoria no atribuye cifras del mercado, prevalencias ni
cuellos de botella específicos a ese documento. La caracterización
cuantitativa de PYMES, latencia real de conectividad y tiempos actuales
de recolección permanece pendiente de levantamiento de campo.

\textbf{Pregunta de trabajo.} ¿Cómo se comporta el flujo autenticado de
gestión de recolecciones en el API Rust+CouchDB bajo las condiciones de
red de laboratorio registradas, y cómo se conservan y sincronizan
cambios locales del cliente cuando se recupera la conectividad?

\section{5. Justificación, alcance y
delimitación}\label{justificaciuxf3n-alcance-y-delimitaciuxf3n}

La solución integra un flujo auditable con persistencia documental y una
cola local, elementos pertinentes para procesos en los que la
conectividad no puede darse por continua. Los incentivos están asociados
a pesajes verificados y sus movimientos se registran en un ledger, con
validaciones de saldo e idempotencia en el canje.

El alcance comprende aplicación web, API, CouchDB, sincronización
PouchDB, pruebas automatizadas y herramientas de carga/telemetría. No
comprende despliegue productivo, evaluación del impacto ambiental,
medición de cambio de conducta de comercios, ni una comparación causal
PostgreSQL-versus-CouchDB. Tampoco se dispone de un servidor PostgreSQL
que replique autenticación, reglas, orquestación y contratos HTTP del
servicio Rust.

\section{6. Objetivos e hipótesis}\label{objetivos-e-hipuxf3tesis}

\textbf{Objetivo general.} Implementar y evaluar en laboratorio un
sistema web de gestión de recolecciones incentivadas para comercios, con
backend Rust, persistencia CouchDB y cliente con capacidad
offline-first.

\textbf{Objetivos específicos.} (a) Implementar el flujo por roles desde
la solicitud hasta el cierre y el ledger de puntos; (b) conservar
localmente datos y cambios pendientes en el cliente y sincronizarlos con
la API; (c) instrumentar y ejecutar corridas repetibles de carga bajo
perfiles de red; y (d) describir estadísticamente las observaciones
válidas sin extrapolar más allá del diseño.

\textbf{Hipótesis experimental.} H0: la media de throughput de corridas
válidas del API no difiere entre las condiciones ideal e intermitente
observadas. H1: al menos una media difiere. El contraste es
exploratorio: hay tres corridas válidas por grupo, las condiciones se
ejecutaron en secuencia y el registro no documenta aleatorización del
orden. Por tanto, la significancia estadística no elimina posibles
efectos temporales o de entorno.

\section{7. Estado del arte y brecha de
investigación}\label{estado-del-arte-y-brecha-de-investigaciuxf3n}

La matriz de antecedentes contiene seis fuentes académicas/técnicas
preseleccionadas y una fuente gris contextual. Dal'Pizzol y Rossetto
(2018) describen una PWA offline-first con PouchDB/CouchDB; Kleppmann et
al.~(2019) exponen principios local-first; Zhang et al.~(2022) estudian
rendimiento de Rust en programas con funcionalidad equivalente; y
trabajos de Li et al.~(2021), Bestari et al.~(2026) y Dillman-Hasso
(2026) abordan incentivos o su persistencia en conductas proambientales.
La matriz separa el nivel de revisión y la aplicabilidad de cada fuente.

El registro consigna cadenas exploratorias ejecutadas mediante un índice
web y consulta de páginas fuente. ScienceDirect se revisó parcialmente a
nivel de registros; no consta búsqueda directa y exportación en IEEE
Xplore, ACM Digital Library o Scopus. No hay conteos reproducibles de
resultados, duplicados y exclusiones ni cribado completo. La brecha es
provisional: en la búsqueda realizada no se identificó un estudio que
reúna comercios salvadoreños, recompensas, sincronización
PouchDB/CouchDB y API Rust bajo red degradada. No se presenta como
ausencia absoluta de literatura ni como revisión sistemática concluida.

\section{8. Marco conceptual}\label{marco-conceptual}

\textbf{Offline-first y local-first.} El cliente conserva datos en el
dispositivo y puede registrar cambios sin conexión; al recuperar
conectividad intenta sincronizarlos. Esto mejora la continuidad
operativa, pero introduce decisiones de conflicto, retención, separación
por identidad y protección del almacenamiento local (Dal'Pizzol \&
Rossetto, 2018; Kleppmann et al., 2019).

\textbf{Trazabilidad e idempotencia.} Cada etapa del proceso tiene actor
y estado; los identificadores deterministas de movimientos previenen
dobles acreditaciones o débitos en reintentos. La verificación del
pesaje es la condición para acreditar puntos.

\textbf{Medición de desempeño.} Una comparación interpretable requiere
describir la unidad medida y mantener equivalentes las cargas. Los
benchmarks de componentes no permiten concluir superioridad de sistemas
completos si difieren las capas y responsabilidades (Zhang et al.,
2022). Por ello el baseline SQL se informa separadamente del flujo HTTP
autenticado.

\section{9. Diseño metodológico}\label{diseuxf1o-metodoluxf3gico}

El trabajo combina desarrollo de software y un experimento de
laboratorio de medidas repetidas independientes por corrida. La
condición de red es el factor de comparación; el throughput HTTP del
flujo k6 es la variable de respuesta primaria del análisis archivado.
Las corridas con éxito de negocio inferior a 100\% o fallos HTTP
distintos de cero se excluyeron del contraste de throughput, pero se
conservaron para el resumen descriptivo de fiabilidad.

La comparación estadística disponible incluye condiciones
\texttt{A1\_ideal\_vus10\_1m} y \texttt{A1\_intermitente\_vus10\_1m},
con tres corridas válidas cada una. El analizador aplicó Shapiro-Wilk,
ANOVA de una vía cuando ambas muestras no rechazaron normalidad,
contraste post hoc Welch, ajuste de Holm, delta de Cliff e intervalo
bootstrap de diferencia de medianas. El mínimo de tres réplicas permite
correr el análisis configurado, pero limita la potencia y la
generalización. No se documentan aleatorización del orden ni replicación
en varios hosts.

Docker Compose proporciona un entorno de laboratorio reproducible con
API, CouchDB, Prometheus, cAdvisor, k6, simulador offline y perfiles de
red \texttt{netem}. Las pruebas de integración y recompensas verifican
rutas y reglas del flujo; no sustituyen el protocolo experimental ni
miden adopción o impacto ambiental.

\section{10. Arquitectura y diseño del
sistema}\label{arquitectura-y-diseuxf1o-del-sistema}

El backend está implementado exclusivamente en Rust y expone una API
HTTP autenticada. CouchDB conserva usuarios, comercios, solicitudes,
pesajes, movimientos de puntos, recompensas y saldos. Se definen índices
Mango para listados por tipo, recompensas activas, movimientos por
comercio, estados/solicitudes de pesajes y ventanas de captura. La
documentación pide verificar el plan mediante \texttt{/\_explain}.

El frontend está desarrollado en TypeScript. PouchDB mantiene una base
local en el navegador y la sincronización utiliza la API Rust como
intermediario de CouchDB. Los documentos guardados pueden llevar estados
locales \texttt{pending}, \texttt{synced}, \texttt{conflict} o
\texttt{rejected}. El simulador \texttt{client-sim} permite escribir
documentos sin HTTP, reenviarlos al reconectar y comprobar su lectura
remota.

El baseline separado usa PostgreSQL y \texttt{pgbench} para hacer upsert
y lectura indexada sobre \texttt{baseline\_documents}. No es un servidor
API alternativo: no reproduce autenticación, autorización, reglas de
negocio, idempotencia, orquestación Rust ni el costo del cliente HTTP.
No se combinan sus observaciones con k6 ni se calcula una comparación
integral entre arquitecturas.

\section{11. Implementación funcional y
seguridad}\label{implementaciuxf3n-funcional-y-seguridad}

Los roles del flujo son \texttt{admin\_plataforma},
\texttt{operador\_acopio}, \texttt{recolector} y \texttt{comercio}. El
comercio crea solicitudes; operador o administración pueden
aprobar/rechazar, asignar y verificar; el recolector registra pesajes de
sus solicitudes; y la recolección se completa cuando existe un pesaje
verificado. Las validaciones incluyen materiales permitidos, límites de
peso/fecha, relación con solicitud y recolector, prevención de
duplicados y transiciones permitidas.

El backend calcula puntos tras la verificación y conserva movimientos
idempotentes. El catálogo permite al comercio consultar recompensas y
solicitar canjes; el servidor comprueba el saldo autoritativo, descuenta
puntos y registra el débito con clave de idempotencia.

La auditoría local encontró el JWT y los metadatos de sesión en
\texttt{localStorage}; datos de negocio y sincronización en PouchDB
sobre IndexedDB. El código de aplicación no cifra esos datos en reposo,
no solicita almacenamiento persistente al navegador y no elimina la base
offline al cerrar sesión. Además, las cuentas \texttt{admin\_plataforma}
usan el mismo nombre de base local. Estas son limitaciones de privacidad
y separación de datos, no garantías de protección. La evaluación
detallada está en \texttt{docs/PRIVACIDAD\_Y\_ALMACENAMIENTO\_LOCAL.md}.

\section{12. Instrumentación, corridas y
metrología}\label{instrumentaciuxf3n-corridas-y-metrologuxeda}

El backend expone \texttt{/metrics} en formato Prometheus con contador
de solicitudes, histograma de latencia por ruta normalizada y contadores
de eventos de negocio. Prometheus está configurado para scrapear
\texttt{api:8080}; cAdvisor expone series de recursos del
host/contendedores. k6 recorre el ciclo autenticado de solicitud,
asignación, pesaje, verificación y puntos.
\texttt{scripts/run\_experiment.ps1} captura JSON por corrida y
actualiza el CSV; el modo offline del simulador reporta convergencia
cuando la sincronización concluye.

La revisión de resultados detectó que los JSON archivados no tienen
muestras válidas de CPU/RAM atribuidas al servicio API
(\texttt{data.result} vacío en esas consultas). Las series agregadas de
cAdvisor no identifican el servicio API y no sirven para atribuirle
consumo. El CSV no incluye CPU ni RAM. No se recolectó ni registró un
conteo directo de hilos; no se informa ni estima esa cantidad. La
metrología de recursos del backend queda pendiente de corregir y validar
en nuevas corridas.

\section{13. Resultados}\label{resultados}

Los valores siguientes son los incluidos en
\texttt{results/analysis.json}. Se muestran con redondeo decimal para
legibilidad; el archivo JSON es la fuente de precisión completa.

{\def\LTcaptype{none} % do not increment counter
\begin{longtable}[]{@{}
  >{\raggedright\arraybackslash}p{(\linewidth - 8\tabcolsep) * \real{0.2000}}
  >{\raggedleft\arraybackslash}p{(\linewidth - 8\tabcolsep) * \real{0.2000}}
  >{\raggedleft\arraybackslash}p{(\linewidth - 8\tabcolsep) * \real{0.2000}}
  >{\raggedleft\arraybackslash}p{(\linewidth - 8\tabcolsep) * \real{0.2000}}
  >{\raggedleft\arraybackslash}p{(\linewidth - 8\tabcolsep) * \real{0.2000}}@{}}
\toprule\noalign{}
\begin{minipage}[b]{\linewidth}\raggedright
Condición (corridas válidas)
\end{minipage} & \begin{minipage}[b]{\linewidth}\raggedleft
Media RPS
\end{minipage} & \begin{minipage}[b]{\linewidth}\raggedleft
Mediana RPS
\end{minipage} & \begin{minipage}[b]{\linewidth}\raggedleft
DE
\end{minipage} & \begin{minipage}[b]{\linewidth}\raggedleft
Mínimo-máximo RPS
\end{minipage} \\
\midrule\noalign{}
\endhead
\bottomrule\noalign{}
\endlastfoot
Ideal, 10 VUs, 1 min (\emph{n} = 3) & 38.9878 & 40.1732 & 4.1772 &
34.3459-42.4442 \\
Intermitente, 10 VUs, 1 min (\emph{n} = 3) & 9.2737 & 9.4904 & 0.5474 &
8.6511-9.6795 \\
\end{longtable}
}

Shapiro-Wilk: condición ideal, \emph{W} = 0.9396, \emph{p} = 0.5258;
intermitente, \emph{W} = 0.8825, \emph{p} = 0.3317. ANOVA de una vía:
\emph{F} = 149.2368, \emph{p} = 0.0002578. Contraste post hoc Welch:
\emph{t} = 12.2163, \emph{p} sin ajustar = 0.0058430, \emph{p} de Holm =
0.0058430; delta de Cliff = 1.0; intervalo bootstrap del 95\% para la
diferencia de medianas = {[}24.6664, 33.7930{]} RPS. Con alfa = 0.05, el
JSON marca rechazo de igualdad para esta comparación.

En el resumen de fiabilidad de todas las corridas, la condición ideal
reunió cinco ejecuciones: tres sin errores y dos excluidas del contraste
de throughput; su éxito medio de flujo fue 0.6833 y la tasa media de
fallos HTTP fue 0.1286. La condición intermitente reunió tres
ejecuciones, todas sin errores; éxito medio de flujo 1.0 y tasa de
fallos HTTP 0.0. La diferencia entre ese resumen y el conjunto
inferencial se debe a que las corridas fallidas no entran al contraste
de throughput.

Estos resultados indican una diferencia estadística de throughput entre
las dos condiciones observadas para el API en este laboratorio. No
estiman una mejora porcentual, no comparan PostgreSQL con CouchDB y no
prueban superioridad de una arquitectura completa. La interpretación
debe considerar el pequeño número de réplicas y la ausencia de
aleatorización documentada.

\section{14. Discusión, viabilidad y
limitaciones}\label{discusiuxf3n-viabilidad-y-limitaciones}

El sistema cubre el flujo funcional y ofrece mecanismos para
sincronización y trazabilidad, mientras que los resultados cuantitativos
actuales se restringen a carga de API en dos perfiles de red. Los
estudios sobre incentivos respaldan medir conducta observada y
persistencia, pero los datos de pruebas técnicas no evidencian que los
comercios aumenten el reciclaje ni que mantengan el comportamiento
después de retirar una recompensa (Bestari et al., 2026; Dillman-Hasso,
2026; Li et al., 2021).

La viabilidad técnica para laboratorio está respaldada por construcción
del stack y pruebas automatizadas documentadas. La viabilidad operativa,
económica, legal y ambiental requiere validación con comercios,
recolectores y entidades locales. Antes de cualquier uso real, las
prioridades son: mejorar protección de token/datos locales, establecer
retención y borrado, aislar bases por identidad, instrumentar CPU/RAM
del contenedor API y conteo verificable de hilos si se mantiene como
métrica, ampliar réplicas y documentar orden aleatorio, y completar una
búsqueda bibliográfica con cribado y conteos reproducibles.

La referencia PostgreSQL/\texttt{pgbench} es útil para caracterizar una
carga SQL de persistencia, pero no es equivalente al flujo autenticado
de k6. Para comparar servicios completos se requiere otro API que
replique contratos, autenticación, autorización y reglas con recursos,
datos, red y cliente equivalentes; hasta entonces no se atribuyen
diferencias a los motores de persistencia.

\section{15. Conclusiones y trabajo
futuro}\label{conclusiones-y-trabajo-futuro}

El Proyecto 15 dispone de un flujo de producto implementado que vincula
comercios, solicitudes, recolectores, pesajes verificados, puntos y
canjes idempotentes, junto con cliente PouchDB y API Rust/CouchDB. El
stack incluye automatización de pruebas, simulación de red y carga.

La evidencia experimental vigente permite una comparación estadística
limitada entre las condiciones ideal e intermitente del mismo API. El
análisis registró diferencias de throughput con Holm \emph{p} =
0.0058430; no hay base para conclusiones comparativas entre APIs
completas ni para medir CPU/RAM del API en las corridas guardadas. El
uso de hilos permanece sin medir.

La revisión bibliográfica es exploratoria y la brecha propuesta es
provisional. La memoria hace visibles esas limitaciones y no sustituye
trabajo de campo ni cierre de protocolo. El trabajo futuro debe
completar la búsqueda sistemática, obtener datos locales primarios,
corregir la atribución de métricas de recursos, ejecutar más réplicas
con orden aleatorio, formalizar privacidad y retención, y evaluar
adopción y persistencia del efecto de los incentivos.

\section{Referencias}\label{referencias}

\begin{hangparas}{0.5in}{1}

Bestari, L. R., Goto, D., & Tran, D. (2026). The effects of monetary incentives and environmental information nudges on household waste sorting behavior: Experimental evidence from community-based waste banks in Indonesia. *Resources, Conservation and Recycling, 232*, Article 108991. https://doi.org/10.1016/j.resconrec.2026.108991

Dal’Pizzol, M., & Rossetto, A. G. de M. (2018). *Offline-first multidatabase Progressive Web Apps com PouchDB/CouchDB*. I Congresso de Tecnologia da Informação do IFSUL Campus Passo Fundo. https://congressoti.passofundo.ifsul.edu.br/flipbook/anais/2018/pdf/anais.pdf

Dillman-Hasso, N. H. (2026). A meta-analysis on the persistent effects of incentives on pro-environmental behavior change interventions. *Journal of Environmental Psychology, 109*, Article 102860. https://doi.org/10.1016/j.jenvp.2025.102860

Kleppmann, M., Wiggins, A., van Hardenberg, P., & McGranaghan, M. (2019). Local-first software: You own your data, in spite of the cloud. *Proceedings of Onward! 2019*, 154-178. https://doi.org/10.1145/3359591.3359737

Li, C., Wang, Y., Li, Y., Huang, Y., & Harder, M. K. (2021). The incentives may not be the incentive: A field experiment in recycling of residential food waste. *Resources, Conservation and Recycling, 168*, Article 105316. https://doi.org/10.1016/j.resconrec.2020.105316

Ministerio de Medio Ambiente y Recursos Naturales. (2023). *Diagnóstico del mercado de reciclaje de El Salvador y articulación de la cadena de valor en el marco de la economía circular*. Biblioteca Ambiental MARN. https://cidoc.ambiente.gob.sv/documento/diagnostico-del-mercado-de-reciclaje-de-el-salvador-y-articulacion-de-la-cadena-de-valor-en-el-marco-de-la-economia-circular/

Mozilla. (s. f.). *Storage quotas and eviction criteria*. MDN Web Docs. https://developer.mozilla.org/en-US/docs/Web/API/Storage_API/Storage_quotas_and_eviction_criteria

PouchDB. (s. f.). *Adapters*. https://pouchdb.com/adapters.html

Zhang, Y., Zhang, Y., Portokalidis, G., & Xu, J. (2022). Towards understanding the runtime performance of Rust. In *Proceedings of the 37th IEEE/ACM International Conference on Automated Software Engineering* (Article 140). https://doi.org/10.1145/3551349.3559494

\end{hangparas}

\section{Apéndice A. Evidencia y archivos
reproducibles}\label{apuxe9ndice-a.-evidencia-y-archivos-reproducibles}

\begin{itemize}
\tightlist
\item
  Observaciones originales por corrida:
  \texttt{results/observations.csv}.
\item
  Análisis estadístico y alcance del baseline:
  \texttt{results/analysis.json}.
\item
  Código del análisis: \texttt{analysis/analizar.py}.
\item
  Registro bibliográfico, búsqueda exploratoria, criterios y
  limitaciones:
  \texttt{docs/revision\_bibliografica/Matriz\_antecedentes\_y\_registro\_busqueda\_Proyecto15.xlsx}.
\item
  Diseño de persistencia y alcance de \texttt{pgbench}:
  \texttt{docs/PERSISTENCIA\_Y\_BASELINE.md}.
\item
  Instrumentación, CPU/RAM e hilos:
  \texttt{docs/MODULO8\_METRICAS\_EXPERIMENTOS.md}.
\item
  Privacidad y almacenamiento local:
  \texttt{docs/PRIVACIDAD\_Y\_ALMACENAMIENTO\_LOCAL.md}.
\item
  Funciones del producto y modo offline:
  \texttt{docs/MODULO7\_FLUJOS\_PRODUCTO.md}.
\end{itemize}

\section{Apéndice B. Comandos de reproducción del
análisis}\label{apuxe9ndice-b.-comandos-de-reproducciuxf3n-del-anuxe1lisis}

Desde la raíz del repositorio, con Python y las dependencias NumPy/SciPy
instaladas:

\begin{Shaded}
\begin{Highlighting}[]
\NormalTok{python }\OperatorTok{.}\NormalTok{\textbackslash{}analysis\textbackslash{}analizar}\OperatorTok{.}\FunctionTok{py} \OperatorTok{.}\NormalTok{\textbackslash{}results\textbackslash{}observations}\OperatorTok{.}\FunctionTok{csv} \OperatorTok{{-}{-}}\NormalTok{output }\OperatorTok{.}\NormalTok{\textbackslash{}results\textbackslash{}analysis}\OperatorTok{.}\FunctionTok{json}
\end{Highlighting}
\end{Shaded}

El PDF se exporta desde este Markdown con Pandoc y XeLaTeX. Al
regenerarlo, primero actualice y revise el JSON de análisis; no agregue
valores de corridas aisladas que no estén en el análisis citado.

\end{document}
